\section{Related Work}
\label{sec:related}

\subsection{Knowledge Distillation for Object Detection}

Knowledge distillation, introduced by \cite{hinton2015distilling}, transfers knowledge from a large teacher network to a smaller student network by matching softened output distributions. For object detection, the challenge extends beyond classification to include localization and feature learning.

Chen et al. \cite{chen2017learning} pioneered KD for detection, proposing to distill both classification and bounding box regression outputs. Subsequent work has explored various knowledge sources: FitNets \cite{romero2014fitnets} focus on intermediate feature alignment; attention transfer methods \cite{zagoruyko2016paying} use attention maps as supervision; and localization-aware distillation \cite{dai2021general} specifically targets bounding box quality.

Recent detection KD frameworks often combine multiple distillation branches, integrating feature and attention distillation for spatial consistency, while relation-based distillation captures structural information. However, these methods are typically validated on clean datasets (COCO, ImageNet), leaving behavior under adverse conditions understudied.

\subsection{Knowledge Distillation Under Domain Shift and Adverse Conditions}

The robustness of KD to domain shift has garnered increasing attention. \cite{sun2016deep} analyze feature alignment across domains, while few-shot adaptation methods explore cross-domain knowledge transfer with limited target data \cite{mirzadeh2020improved}.

Adverse weather conditions present a particularly challenging domain shift scenario. \cite{guo2021meets} survey the intersection of KD and image restoration, noting that degradations like fog and haze degrade both teacher and student performance. \cite{du2022learning} propose learning weather-invariant representations through adversarial training, but do not isolate specific failure mechanisms.

Several works have addressed fog-specific object detection. \cite{li2018benchmarking} establish benchmarks for dehazing and visibility restoration. \cite{sakaridis2018semantic} provide the Foggy Cityscapes dataset demonstrating significant performance degradation under synthetic fog. While these works acknowledge that standard detectors struggle in fog, they do not analyze how knowledge transfer specifically breaks down.

\subsection{Mechanism-Oriented Analysis of Knowledge Distillation}

Beyond proposing new methods, a growing body of work seeks to understand \emph{why} KD works or fails. \cite{phuong2019functional} provide theoretical analysis of the function learned by distillation, while \cite{cho2019efficacy} examine which aspects of teacher knowledge actually transfer.

Empirical mechanism studies often focus on the teacher-student relationship. \cite{mirzadeh2020improved} investigate the gap between teacher and student capacity, finding that imperfect teachers can still provide useful supervision. \cite{gou2021knowledge} survey various knowledge types and their transfer mechanisms.

For failure analysis, several hypotheses recur in the literature. \textbf{Distribution mismatch} posits that domain shift disrupts the teacher output distributions that KD relies upon \cite{wang2020cross}. \textbf{Uncertainty amplification} suggests that teacher predictions become unreliable under degraded inputs, amplifying errors in student training \cite{kendall2017uncertainties}. \textbf{Representation misalignment} proposes that feature spaces diverge across domains, impeding feature-based distillation \cite{sun2016deep}.

However, these mechanisms are rarely tested against each other within a unified framework. Our work fills this gap by empirically evaluating multiple hypotheses within a controlled experimental design, determining which mechanisms receive direct support from observed behavior.

% TODO: Additional recent references to be added:
% - [ ] SPIC journal guidelines on detection under degraded conditions
% - [ ] Recent arXiv preprints on KD robustness (2023-2024)
% - [ ] Domain adaptation literature for adverse weather
